
\section{\textbf{\texttt{name\_create()}} ORM Method}

\begin{frame}[fragile]{\textbf{\texttt{name\_create()}} Method}
	\begin{itemize}
		\item The \textbf{\texttt{name\_create()}} method creates a record using only a display name.
		\item It is used by many2one fields when the user types a new value and quick-creates it.
		\item Internally, it usually calls \texttt{create()} with the \texttt{name} field.
\begin{block}{Method Syntax}
\begin{lstlisting}
@api.model
def name_create(self, name):
	return super().name_create(name)
\end{lstlisting}
	\begin{itemize}
		\item Important parts: \textbf{\texttt{@api.model}}, \textbf{\texttt{name}}.
	\end{itemize}
\end{block}
	\end{itemize}
\end{frame}

\begin{frame}[fragile]{Breaking Down the Method Signature}
	\begin{block}{}
		\begin{itemize}
			\item \textbf{\texttt{@api.model}}: Runs at the model level.
			\item \textbf{\texttt{self}}: The model on which the new record is created.
			\item \textbf{\texttt{name}}: A string that becomes the record's name / display name.
\begin{lstlisting}
name = "Ahmed Mohamed"
\end{lstlisting}
			\item No full \texttt{vals} dictionary is required by the caller.
			\item Extra required fields must be handled by defaults or an override.
		\end{itemize}
	\end{block}
\end{frame}

\begin{frame}[fragile]{What Does \texttt{name\_create()} Return?}
	\begin{block}{}
		\begin{itemize}
			\item It returns a tuple: \texttt{(id, display\_name)}.
\begin{lstlisting}
result = self.env['student.student'].name_create('Ahmed Mohamed')
print(result)
\end{lstlisting}
			Output is (example):
\begin{lstlisting}
(29, 'Ahmed Mohamed')
\end{lstlisting}
			\item First value: database ID of the new record.
			\item Second value: display name shown in the UI.
		\end{itemize}
	\end{block}
\end{frame}

\begin{frame}[fragile]{How the UI Uses \texttt{name\_create()}}
	\begin{itemize}
		\item In a many2one field, the user types a name that does not exist.
		\item Odoo offers ``Create \ldots'' / quick create.
		\item Selecting that option calls \texttt{name\_create()}.
\begin{lstlisting}
# Conceptual flow
# typed value: "Grade 12"
self.env['school.class'].name_create('Grade 12')
\end{lstlisting}
		\item This is why many2one quick-create only needs a name.
	\end{itemize}
\end{frame}

\begin{frame}[fragile]{Override Pattern}
	When a model has more required fields than \texttt{name}:
\begin{lstlisting}
@api.model
def name_create(self, name):
	record = self.create({
		'name': name,
		'active': True,
		'gender': 'other',
	})
	return record.id, record.display_name
\end{lstlisting}
	\begin{itemize}
		\item Override \texttt{name\_create()} when quick-create must fill extra values.
	\end{itemize}
\end{frame}

\begin{frame}[fragile]{Method Call}
\begin{block}{Method Call}
\begin{lstlisting}
record_id, display_name = self.env['student.student'].name_create(
	'Ahmed Mohamed'
)
\end{lstlisting}
	\begin{itemize}
		\item Here, \textbf{\texttt{.name\_create(...)}} is the \textbf{method call}.
		\item Return type reminder: \textbf{tuple (id, display\_name)}.
	\end{itemize}
\end{block}
\end{frame}
